% Intended to REPLACE the "Comparison with healthcare and the role of
% gradient noise" paragraph at the end of \subsection{Synthetic
% Multidimensional Knapsack} (sec:exp-knapsack) and the paragraph that
% currently follows it. The existing gradient-cosine figure
% (fig:gradient-knapsack) reference is preserved; only the surrounding
% prose changes to incorporate the mu-sweep and normalized-PCGrad
% findings without changing any numbers in the main knapsack tables.

\textit{Prediction-anchor dependence under stochastic decision gradients.}
A striking feature of the knapsack experiment is that removing the
prediction anchor entirely is not a safe operating point for FDFL. With
$\mu = 0$ in the loss $L_{\mathrm{regret}} + \mu\,L_{\mathrm{pred}} +
\lambda F$, i.e.\ when the predictor is trained on decision regret and
prediction-fairness alone, FDFL diverges on a substantial fraction of
random seeds and attains average normalized regret of $0.84 \pm 0.36$
at the redesigned $n_{\mathrm{train}} = 300$, $\alpha = 2$ MAD cell
(5 seeds, 30 SPSA steps per $\lambda$). This is not a variance effect:
two of the five seeds terminate with the predictor pushed into a
degenerate region where subsequent SPSA estimates provide no useful
gradient. Introducing even a small fixed weight restores stability
completely: FDFL at $\mu = 0.1$ drops to $0.197 \pm 0.004$, at
$\mu = 0.5$ to $0.186 \pm 0.005$, and at $\mu = 1$ (FDFL-Scal) to
$0.175 \pm 0.007$. The choice among these three $\mu$ values is
second-order---they all sit in the same performance band as the
scalarized FDFL-Scal reported in Table~\ref{tab:knapsack-a2}. The
practical message is that a positive prediction weight acts as a cheap
stabilizer in the stochastic-gradient regime, and the empirical
benefit of leaving $\mu = 0$, which the original FDFL specification
allows, depends on the availability of an accurate decision-gradient
estimator.

\textit{MOO in the stochastic-gradient regime.}
The contrasting behavior of the multi-objective methods across the two
experiments is informative. In the healthcare setting at $\alpha = 2$,
FPLG and FDFL-PCGrad jointly achieved the lowest regret (within
$0.0007$ of each other, both below every two-stage baseline); in the
knapsack setting, FDFL-Scal and $\mu$-anchored FDFL variants are on
top and the MOO methods are noticeably less reliable. Two structural
differences explain the shift. First, the \textit{gradient computation}
differs: the healthcare experiment uses the exact analytical Jacobian,
while the knapsack relies on SPSA perturbation-based approximation.
MOO algorithms such as PCGrad depend on accurate gradient directions
to identify and resolve conflicts between objectives; when the decision
gradient is noisy, the conflict-resolution projections can remove
useful signal or introduce spurious corrections. Second, the
\textit{gradient alignment structure} differs (Figure~\ref{fig:gradient-knapsack}):
prediction--fairness cosines and decision--prediction cosines are all
small in the knapsack, indicating near-orthogonal objectives for which
projection-based combination offers no advantage over a scalarized sum.

A natural attempt at mitigation is to equalize the per-objective
gradient magnitudes before PCGrad's conflict-resolution step, since
the raw decision-regret norms under SPSA exceed the prediction-loss
and fairness norms by roughly three orders of magnitude. We implemented
this as an optional per-objective normalization inside the PCGrad
handler, rescaling the summed direction back by the mean of the
original gradient norms to preserve the step size. Normalization
removes the directional imbalance on most seeds, but it does not
eliminate the stochastic-gradient failure mode: over nine seeds on the
redesigned knapsack, eight produce stable runs with regret
$0.189 \pm 0.024$ (competitive with FDFL-Scal), while one seed still
diverges catastrophically ($0.9999$, train regret $1.6 \times 10^9$),
an observed failure rate of $\approx 11\%$. The residual failure mode
is that the rescaling still transmits SPSA magnitude spikes to the
effective step size, and a single oversized step into a degenerate
region of parameter space is not recoverable under subsequent noisy
gradient estimates. On the analytic-gradient healthcare task the same
normalization has a small and $\alpha$-dependent effect: a $\approx 15\%$
reduction in regret relative to the unnormalized variant at
$\alpha = 0.5$ across all four fairness measures, and a $\approx 4\%$
increase at $\alpha = 2$. Appendix~\ref{app:mu-followup} reports the
full results of the $\mu$-sweep and the PCGrad normalization
experiments.

These findings suggest a refined practical guideline. For FDFL training,
(i)~when decision gradients are computed analytically and the objective
gradients exhibit meaningful conflict, MOO-based combination can match
or slightly improve upon scalarized FDFL; (ii)~when decision gradients
are approximate, a fixed prediction weight $\mu > 0$ is the first
stabilizer to add, and it largely closes the gap between FDFL variants;
(iii)~per-objective gradient normalization for PCGrad is a useful
configuration option, but it should be framed as an $\alpha$-dependent
trade-off rather than a universal improvement.
